\documentclass[11pt,a4paper]{article}

\usepackage[utf8]{inputenc}
\usepackage[T1]{fontenc}
\usepackage{amsmath,amssymb,amsthm}
\usepackage{booktabs}
\usepackage{graphicx}
\usepackage{hyperref}
\usepackage[margin=1in]{geometry}
\usepackage{enumitem}
\usepackage{array}
\usepackage{float}
\usepackage{caption}
\usepackage{natbib}
\bibliographystyle{plainnat}

\newtheorem{theorem}{Theorem}
\newtheorem{proposition}[theorem]{Proposition}
\newtheorem{lemma}[theorem]{Lemma}
\newtheorem{corollary}[theorem]{Corollary}
\newtheorem{definition}[theorem]{Definition}
\newtheorem{remark}[theorem]{Remark}

\title{A Universal Cost Function for Musical Intervals:\\
Connecting Harmonic Structure, Emotional Valence,\\
and Chromatic Architecture Through a Single Formula}

\author{Jonathan Washburn\\
Recognition Science Institute\\
\texttt{jon@recognitionscience.com}}

\date{March 2026}

\begin{document}

\maketitle

\begin{abstract}
We study the function $J(r)=\tfrac{1}{2}(r+r^{-1})-1$ as a universal
cost measure for musical frequency ratios.  The accompanying Lean
formalization proves a uniqueness theorem for $J$ on the positive reals
under an explicit hypothesis bundle comprising inversion symmetry, unit
normalization, strict convexity, log-coordinate calibration, continuity
on $(0,\infty)$, a cosh-add functional identity, and the regularity
hypotheses used in the d'Alembert/ODE reduction.  Applying $J$ to
standard just-intonation ratios produces a complete ordering of
musical intervals---unison, minor third, major third, perfect fourth,
tritone, perfect fifth, octave---that differs from traditional
consonance rankings in specific, testable ways: notably, the perfect
fifth has \emph{higher} $J$-cost than the minor third, major third,
and perfect fourth.  We show that the fifth's structural importance
arises instead from its role as the bridge ratio $12/8=3/2$ linking the
chromatic octave to an 8-slot oscillation basis, with the twelve-fifths
walk closing the octave to within $<2\%$ (the Pythagorean comma).  We
introduce a signed asymmetry function $\sigma(r)=r-r^{-1}$ that
serves as a comparative valence proxy: larger $\sigma$ corresponds to
brighter affect, and the major third has larger positive skew than the
minor third.  Finally, we extend $J$ from intervals to chords via a
directed pairwise aggregate, defining coherence as
$\mathrm{Coh}(C)=1/(1+\mathrm{Conflict}(C))$, and prove that lower
aggregate $J$-cost implies higher coherence.  The core theorem surface
cited explicitly in the paper has corresponding machine-checked proofs
in the Lean~4 proof assistant.
\end{abstract}

\paragraph{Keywords.} Consonance, dissonance, cost function, musical
intervals, valence, circle of fifths, Pythagorean comma, formal
verification, just intonation

% ====================================================================
\section{Introduction}
\label{sec:intro}

Why do simple frequency ratios sound consonant?  This question has
occupied acousticians since Pythagoras, yet no single quantitative
measure commands universal agreement.  Helmholtz's roughness model
\citep{helmholtz1863} appeals to beating between partials;
Plomp--Levelt critical-bandwidth theory \citep{plomp1965} grounds
dissonance in the basilar membrane; harmonic-template models
\citep{parncutt2011} invoke pattern recognition in the auditory
cortex.  Each captures part of the phenomenon.  The present paper asks a
different question: how much of interval ordering, bridge structure,
and signed asymmetry can be organized by a single closed-form cost
function with no fitted parameters?

This paper proposes such a function.  The cost
\begin{equation}
  \label{eq:jcost-intro}
  J(r) = \tfrac{1}{2}\!\left(r + r^{-1}\right) - 1
\end{equation}
measures how far a positive ratio $r$ deviates from unity, and it does
so with no free parameters beyond the definition itself.  We show:

\begin{enumerate}[label=(\arabic*)]
  \item the Lean development proves a uniqueness theorem for $J$ on
    $\mathbb{R}_{>0}$ under an explicit symmetry/normalization/
    convexity/calibration/functional-equation hypothesis bundle
    (Section~\ref{sec:axioms});
  \item $J$ produces a complete, strict ordering of the standard
    just-intonation intervals (Section~\ref{sec:hierarchy});
  \item the perfect fifth's structural importance comes from a
    chromatic-bridge property, not from having the lowest cost
    (Section~\ref{sec:fifth});
  \item a signed companion function $\sigma(r)=r-r^{-1}$ gives a
    parameter-free model of emotional valence
    (Section~\ref{sec:valence});
  \item extending $J$ to chords via pairwise aggregation yields a
    monotone conflict-coherence framework
    (Section~\ref{sec:conflict}).
\end{enumerate}

\noindent The core theorem surface cited by name below has been
formally verified in the Lean~4 proof assistant \citep{moura2021} using
the Mathlib library; the verification is summarised in
Section~\ref{sec:lean}.

% ====================================================================
\section{Axiomatic Derivation of the Cost Function}
\label{sec:axioms}

\subsection{Motivating axioms}

We seek a function $F\colon\mathbb{R}_{>0}\to\mathbb{R}$ satisfying:

\begin{enumerate}[label=(A\arabic*)]
  \item \textbf{Inversion symmetry.}  $F(r)=F(r^{-1})$ for all
    $r>0$.
  \item \textbf{Unit normalisation.}  $F(1)=0$.
  \item \textbf{Strict convexity.}  $F$ is strictly convex on
    $(0,\infty)$.
  \item \textbf{Calibration.}  In logarithmic coordinates, define
    $G(t):=F(e^t)$.  We require $G''(0)=1$.
\end{enumerate}

\noindent Axiom~(A1) says that the cost of an interval is the same
whether one hears it as ascending or descending.  Axiom~(A2) says
unison has zero cost.  Axiom~(A3) ensures a unique global minimum at
$r=1$.  Axiom~(A4) fixes the overall scale so that $F$ is not merely
determined up to a positive multiple.

\medskip

\noindent These four conditions are the conceptual package motivating
the paper.  The machine-checked uniqueness theorem in Lean uses a
strengthened explicit hypothesis bundle: in addition to the items above,
it assumes continuity on $(0,\infty)$, a cosh-add functional identity,
and the regularity hypotheses needed for the d'Alembert-to-ODE step.

\subsection{Derivation}

\begin{theorem}[Squared-ratio form]
\label{thm:sq-form}
For any $r>0$,
\[
  J(r) = \frac{(r-1)^2}{2r}.
\]
\end{theorem}

\begin{proof}
Expanding $J(r)=\tfrac{1}{2}(r+r^{-1})-1$ and multiplying through by
$r>0$ gives $rJ(r)=\tfrac{1}{2}(r^2+1)-r=\tfrac{1}{2}(r-1)^2$.
\end{proof}

\noindent The squared-ratio form makes properties (A1)--(A4)
transparent: the numerator $(r-1)^2$ is symmetric under
$r\leftrightarrow r^{-1}$ when divided by $r$, vanishes at $r=1$, and
is strictly convex.

\begin{theorem}[Basic properties of $J$]
\label{thm:basic}
\leavevmode
\begin{enumerate}[label=(\roman*)]
  \item $J(1)=0$.
  \item $J(r)=J(r^{-1})$ for all $r>0$.
  \item $J(r)>0$ for all $r\neq 1$, $r>0$.
  \item $J$ is strictly convex on $(0,\infty)$.
\end{enumerate}
\end{theorem}

\begin{proof}
Properties (i)--(iii) follow immediately from the squared-ratio form:
$(r-1)^2>0$ for $r\neq 1$ and $2r>0$.  Strict convexity (iv) is
verified by computing $J''(r)=r^{-3}>0$ for $r>0$.
\end{proof}

\begin{theorem}[Uniqueness]
\label{thm:unique}
Let $F\colon\mathbb{R}_{>0}\to\mathbb{R}$ be continuous, satisfy
(A1)--(A4), satisfy the cosh-add identity
\[
  G(s+t)+G(s-t)=2G(s)G(t)+2G(s)+2G(t), \qquad G(t):=F(e^t),
\]
and satisfy the regularity hypotheses needed to identify the associated
continuous d'Alembert solution with $\cosh$.
Then $F=J$ on $(0,\infty)$.
\end{theorem}

\begin{proof}[Proof sketch]
Define $H(t):=G(t)+1$.  Then $H$ satisfies the classical d'Alembert
equation $H(s+t)+H(s-t)=2H(s)H(t)$, with $H(0)=1$, $H$ continuous and
even (from A1), and $H''(0)=G''(0)=1$ (from A4).  The only continuous
solution of d'Alembert's equation with these boundary conditions is
$H(t)=\cosh t$ \citep{aczel1966}.  Hence $G(t)=\cosh t-1$, and
unwinding $F(r)=G(\log r)$ gives $F(r)=\cosh(\log r)-1=J(r)$.
\end{proof}

\begin{remark}
The full formal proof of Theorem~\ref{thm:unique}, including the
reduction from d'Alembert's equation to the cosh solution, is available
through the theorems \texttt{CostUniqueness.T5\_uniqueness\_complete}
and \texttt{CostUniqueness.unique\_cost\_on\_pos}.  In the latter, the
complete machine-checked premise set is packaged into the structure
\texttt{UniqueCostAxioms}.  Several regularity hypotheses for the ODE
step are carried as explicit premises; they do not involve any
\texttt{sorry} or \texttt{admit}.
\end{remark}

\subsection{The cosh representation}

The identity $J(r)=\cosh(\log r)-1$ connects $J$ directly to the
hyperbolic cosine.  Because $\cosh$ is the unique even solution of
$y''=y$ with $y(0)=1$, this representation makes the
information-theoretic content of $J$ visible: the cost of a ratio $r$
is the excess of the geometric mean of $r$ and $r^{-1}$ over unity.

% ====================================================================
\section{The Interval-Cost Hierarchy}
\label{sec:hierarchy}

\subsection{\texorpdfstring{$J$-cost}{J-cost} of standard intervals}

Applying $J$ to the standard just-intonation ratios yields exact
rational values:

\begin{table}[H]
\centering
\caption{$J$-cost of standard musical intervals in just intonation.}
\label{tab:jcost}
\begin{tabular}{@{}llcl@{}}
\toprule
Interval & Ratio $r$ & $J(r)$ & Decimal \\
\midrule
Unison        & $1/1$    & $0$          & $0$ \\
Minor third   & $6/5$    & $1/60$       & $0.0167$ \\
Major third   & $5/4$    & $1/40$       & $0.0250$ \\
Perfect fourth& $4/3$    & $1/24$       & $0.0417$ \\
Tritone       & $45/32$  & $169/2880$   & $0.0587$ \\
Perfect fifth & $3/2$    & $1/12$       & $0.0833$ \\
Octave        & $2/1$    & $1/4$        & $0.2500$ \\
\bottomrule
\end{tabular}
\end{table}

\begin{theorem}[Complete interval-cost ordering]
\label{thm:hierarchy}
\[
  J(1) \;<\; J\!\left(\tfrac{6}{5}\right)
       \;<\; J\!\left(\tfrac{5}{4}\right)
       \;<\; J\!\left(\tfrac{4}{3}\right)
       \;<\; J\!\left(\tfrac{45}{32}\right)
       \;<\; J\!\left(\tfrac{3}{2}\right)
       \;<\; J(2).
\]
\end{theorem}

\begin{proof}
Each inequality reduces to a comparison of explicit rationals.  The
formal proof is the Lean theorem
\texttt{Consonance.extended\_ratio\_cost\_hierarchy}.
\end{proof}

\subsection{Superparticular intervals}

A superparticular ratio is $(n+1)/n$ for a positive integer $n$.  These
are the building blocks of just intonation: the octave ($n=1$), fifth
($n=2$), fourth ($n=3$), major third ($n=4$), and minor third ($n=5$).

\begin{theorem}[Superparticular cost formula]
\label{thm:super}
For every positive integer $n$,
\[
  J\!\left(\frac{n+1}{n}\right) = \frac{1}{2n(n+1)}.
\]
\end{theorem}

\begin{proof}
Direct computation from the squared-ratio form:
$J\bigl(\frac{n+1}{n}\bigr) = \frac{(1/n)^2}{2\cdot(n+1)/n}
= \frac{1}{2n(n+1)}$.
\end{proof}

\begin{corollary}[Superparticular cost is strictly decreasing]
\label{cor:super-dec}
If $1\leq m<n$, then
$J\bigl(\frac{n+1}{n}\bigr) < J\bigl(\frac{m+1}{m}\bigr)$.
\end{corollary}

\begin{proof}
$2n(n+1)>2m(m+1)$ when $n>m\geq 1$, so $1/(2n(n+1))<1/(2m(m+1))$.
\end{proof}

\noindent This means that \emph{smaller} superparticular intervals are
\emph{less costly in the present model}---the minor third ($1/60$) beats the major third
($1/40$), which beats the fourth ($1/24$), which beats the fifth
($1/12$), which beats the octave ($1/4$).  The ordering is the reverse
of what traditional consonance rankings predict for the fifth versus the
thirds, and this reversal is a distinguishing, testable prediction of
the $J$-cost framework.

\subsection{Comparison with existing consonance theories}

\begin{table}[H]
\centering
\caption{Ranking of intervals by three frameworks.  Lower rank = more
consonant.  The $J$-cost ranking diverges at the fifth.}
\label{tab:compare}
\begin{tabular}{@{}lccc@{}}
\toprule
Interval & $J$-cost rank & Helmholtz rank & Harmonic template rank \\
\midrule
Unison        & 1 & 1 & 1 \\
Minor third   & 2 & 5 & 5 \\
Major third   & 3 & 4 & 4 \\
Perfect fourth& 4 & 3 & 3 \\
Perfect fifth & 6 & 2 & 2 \\
Octave        & 7 & 6 & 6 \\
\bottomrule
\end{tabular}
\end{table}

\noindent The $J$-cost ordering agrees with tradition on unison (most
consonant) and octave (least consonant among the standard intervals),
but disagrees on the relative position of the fifth.  In the
Helmholtz and harmonic-template frameworks, the perfect fifth is the
second-most consonant interval.  In the $J$-cost framework, the fifth
ranks sixth---above only the octave.  This is because $J$ measures
deviation from unity: $3/2$ deviates more than $6/5$ or $5/4$.

The disagreement is not a defect but a \emph{model-side prediction}.
Within the present framework, $J$ should be read as an interval-strain
ordering rather than as a complete replacement for all psychoacoustic
preference judgments.  This is exactly where the $J$-cost picture
becomes empirically distinctive.

% ====================================================================
\section{The Structural Role of the Perfect Fifth}
\label{sec:fifth}

If the fifth is not the lowest-cost interval, why is it so
structurally important in Western music?  We argue that the answer
lies in its role as a \emph{bridge ratio} between two independently
motivated counting systems.

\subsection{The 12/8 bridge}

\begin{definition}[Chromatic bridge ratio]
\label{def:bridge}
The \emph{chromatic bridge ratio} is the quotient of the number of
semitones per octave (12) and the number of slots in the DFT-8
oscillation basis (8):
\[
  R_{\mathrm{bridge}} = \frac{12}{8} = \frac{3}{2}.
\]
\end{definition}

\noindent The fact that this ratio \emph{is} the perfect fifth is
the key structural observation.  It means that the perfect fifth is
the interval that connects the chromatic counting system (12 semitones)
to the oscillation counting system (8 time-slots).

\subsection{The Pythagorean comma as a closure witness}

Walking twelve perfect fifths around the circle of fifths and reducing
by seven octaves gives:
\begin{equation}
  \label{eq:comma}
  \frac{(3/2)^{12}}{2^7} = \frac{3^{12}}{2^{19}} \approx 1.01364.
\end{equation}

\begin{theorem}[Pythagorean comma bounds]
\label{thm:comma}
\[
  1 \;<\; \frac{(3/2)^{12}}{2^7} \;<\; 1.02.
\]
\end{theorem}

\begin{proof}
Both inequalities are verified by exact rational arithmetic.  The
formal proofs are \texttt{CircleOfFifths.pythagorean\_comma\_positive}
and \texttt{CircleOfFifths.pythagorean\_comma\_small}.
\end{proof}

\noindent The twelve-fifths walk closes the octave with less than 2\%
error.  This makes the fifth a \emph{low-error closure witness} for
the chromatic system---a structural bridge, not a cost optimum.

\begin{theorem}[The fifth is structural, not lowest-cost]
\label{thm:fifth-not-lowest}
\[
  J\!\left(\tfrac{6}{5}\right) < J\!\left(\tfrac{3}{2}\right),\quad
  J\!\left(\tfrac{5}{4}\right) < J\!\left(\tfrac{3}{2}\right),\quad
  J\!\left(\tfrac{4}{3}\right) < J\!\left(\tfrac{3}{2}\right).
\]
\end{theorem}

\noindent The formal proof is
\texttt{ChromaticDerivation.fifth\_is\_structural\_not\_lowest\_cost}.

\subsection{Harmonic relationships}

The product rule $\text{fifth}\times\text{fourth}=\text{octave}$, i.e.\
$(3/2)\cdot(4/3)=2$, is verified formally as
\texttt{HarmonicModes.fifth\_times\_fourth\_eq\_octave}.  This reflects
the well-known complementarity of the fifth and fourth within the
octave and follows purely from arithmetic of the superparticular
ratios.

% ====================================================================
\section{Emotional Valence as Ratio Asymmetry}
\label{sec:valence}

The cost function $J$ is symmetric under inversion: $J(r)=J(r^{-1})$.
But the \emph{experience} of an interval is not symmetric---a major
third sounds ``brighter'' than a minor third.  We capture this
asymmetry with a signed companion function.

\subsection{The ledger-skew function}

\begin{definition}[Ledger skew]
\label{def:skew}
For $r>0$, the \emph{ledger skew} is
\begin{equation}
  \label{eq:skew}
  \sigma(r) = r - r^{-1}.
\end{equation}
\end{definition}

\noindent The function $\sigma$ captures the antisymmetric part of the
ratio, complementing the symmetric part captured by $J$.  Together,
$J$ and $\sigma$ decompose the ratio into a ``magnitude of deviation''
and a ``direction of deviation.''

\begin{theorem}[Properties of $\sigma$]
\label{thm:skew-props}
\leavevmode
\begin{enumerate}[label=(\roman*)]
  \item $\sigma(1)=0$.
  \item $\sigma(r^{-1})=-\sigma(r)$ for all $r\neq 0$
    \textup{(}antisymmetry\textup{)}.
  \item $\sigma(r)>0$ for all $r>1$.
  \item For $1<r_1<r_2<2$, $\sigma(r_1)<\sigma(r_2)$
    \textup{(}skew increases with interval size within the
    octave\textup{)}.
\end{enumerate}
\end{theorem}

\begin{proof}
Properties (i)--(iii) are immediate.  Property (iv): on $(1,2)$ the
derivative $\sigma'(r)=1+r^{-2}>0$, so $\sigma$ is strictly increasing.
Formally, the Lean proof
\texttt{Valence.skew\_increases\_with\_interval\_size} verifies (iv)
using the bound $r_2^{-1}<r_1^{-1}$ when $r_1<r_2$.
\end{proof}

\subsection{Major versus minor}

\begin{theorem}[Major-minor brightness distinction]
\label{thm:major-minor}
\[
  0
  \;<\;
  \sigma\!\left(\tfrac{6}{5}\right)=\tfrac{11}{30}
  \;<\;
  \sigma\!\left(\tfrac{5}{4}\right)=\tfrac{9}{20}.
\]
\end{theorem}

\begin{proof}
Direct computation.  The formal proofs are
\texttt{Valence.major\_third\_skew}, \texttt{Valence.minor\_third\_skew},
\texttt{Valence.major\_third\_skew\_pos},
\texttt{Valence.minor\_third\_skew\_pos}, and
\texttt{Valence.major\_skew\_gt\_minor\_skew}.
\end{proof}

\noindent The major third has higher ledger skew than the minor third.
In the present framework this supports a \emph{comparative brightness
ordering}: both thirds have positive skew, but the major third is more
``expansive'' (further from reciprocal symmetry on the positive side)
than the minor third.  The quantitative skew difference is:

\begin{theorem}[Valence difference]
\label{thm:val-diff}
\[
  \sigma\!\left(\tfrac{5}{4}\right) - \sigma\!\left(\tfrac{6}{5}\right)
  = \frac{1}{12}.
\]
\end{theorem}

\begin{proof}
The formal proof is
\texttt{Valence.valence\_difference\_one\_twelfth}.
\end{proof}

\begin{remark}
The valence model predicts that any interval with $r>1$ has positive
skew, while its inversion $r^{-1}<1$ has negative skew.  The major/minor
comparison is therefore not a sign flip but a difference in the
magnitude of positive skew.  Formally, the relevant Lean ingredients are
\texttt{ledgerSkew\_pos\_above\_one},
\texttt{ledgerSkew\_neg\_below\_one},
\texttt{ledgerSkew\_antisymmetric}, and
\texttt{skew\_increases\_with\_interval\_size}.
\end{remark}

% ====================================================================
\section{From Intervals to Chords: Conflict and Coherence}
\label{sec:conflict}

\subsection{Chord conflict}

\begin{definition}[Chord conflict]
\label{def:conflict}
For a chord $C=\{f_1,\ldots,f_n\}$ with $f_i>0$, the
\emph{chord conflict} is the directed aggregate of pairwise $J$-costs:
\begin{equation}
  \label{eq:conflict}
  \mathrm{Conflict}(C) = \sum_{i\neq j} J\!\left(\frac{f_i}{f_j}\right).
\end{equation}
\end{definition}

\noindent Because $J(r)=J(r^{-1})$, the directed aggregate equals
twice the unordered pair sum $\sum_{i<j}J(f_i/f_j)$ when all notes
are distinct.  We use the directed form because it matches the
companion Lean implementation exactly.

\begin{theorem}[Conflict is nonneg]
\label{thm:conflict-nonneg}
$\mathrm{Conflict}(C)\geq 0$ for every chord $C$.
\end{theorem}

\begin{proof}
Each summand $J(f_i/f_j)\geq 0$ by Theorem~\ref{thm:basic}(iii),
and the diagonal terms ($i=j$) contribute $J(1)=0$.  The formal proof
is \texttt{ModeConflict.chordConflict\_nonneg}.
\end{proof}

\subsection{Chord coherence}

\begin{definition}[Chord coherence]
\label{def:coherence}
\begin{equation}
  \label{eq:coherence}
  \mathrm{Coh}(C) = \frac{1}{1+\mathrm{Conflict}(C)}.
\end{equation}
\end{definition}

\begin{proposition}[Coherence bounds]
\label{prop:coh-bounds}
$0<\mathrm{Coh}(C)\leq 1$ for every chord $C$.
\end{proposition}

\begin{proof}
Since $\mathrm{Conflict}(C)\geq 0$, the denominator is $\geq 1$, so
$\mathrm{Coh}\leq 1$.  Positivity of the denominator gives
$\mathrm{Coh}>0$.  Formal proofs:
\texttt{ModeConflict.chordCoherence\_pos} and
\texttt{ModeConflict.chordCoherence\_le\_one}.
\end{proof}

\begin{theorem}[Monotonicity]
\label{thm:mono}
If $\mathrm{Conflict}(C_1)\leq\mathrm{Conflict}(C_2)$, then
$\mathrm{Coh}(C_2)\leq\mathrm{Coh}(C_1)$.
\end{theorem}

\begin{proof}
The map $x\mapsto 1/(1+x)$ is decreasing on $[0,\infty)$.  Formal
proof: \texttt{ModeConflict.lowerConflict\_higherCoherence}.
\end{proof}

\noindent Theorem~\ref{thm:mono} is the central structural result
for multi-note applications: among chords matched for other
properties, those with lower aggregate $J$-cost are more coherent.
This gives a parameter-free objective function for algorithmic
composition and for downstream application layers that use coherence as
an alignment modifier.

\subsection{Effective alignment}

In a downstream application layer \citep{washburn2026envelope}, the
coherence score modulates a base alignment to yield an effective
alignment $a_{\mathrm{eff}}=a\cdot\mathrm{Coh}(C)$.  The monotonicity
theorem above guarantees that lower-conflict chords preserve at least as
much alignment as higher-conflict chords, all else equal.  The
corresponding downstream Lean theorem is
\texttt{SoundHealing.lowerConflict\_reduces\_strain\_more}.

% ====================================================================
\section{Formal Verification}
\label{sec:lean}

The core theorem surface cited in this paper has corresponding
machine-checked proofs in the Lean~4 proof assistant \citep{moura2021}.
The cited modules build with zero \texttt{sorry} and zero
\texttt{admit}.

\begin{table}[H]
\centering
\small
\caption{Core Lean~4 theorem surface used in the paper.}
\label{tab:lean}
\begin{tabular}{@{}ll@{}}
\toprule
Claim cluster & Key theorem(s) \\
\midrule
\texttt{Basic cost facts}
  & \texttt{Jcost\_unit0} \\
  & \texttt{Jcost\_symm} \\
  & \texttt{Jcost\_nonneg} \\
  & \texttt{Jcost\_pos\_of\_ne\_one} \\
  & \texttt{Jcost\_eq\_sq} \\
  & \texttt{Jcost\_eq\_zero\_iff} \\
\texttt{Convexity}
  & \texttt{Jcost\_strictConvexOn\_pos} \\
\texttt{Uniqueness}
  & \texttt{T5\_uniqueness\_complete} \\
  & \texttt{unique\_cost\_on\_pos} \\
\texttt{Interval values}
  & \texttt{fifth\_jcost}, \texttt{fourth\_jcost}, etc. \\
  & \texttt{fifth\_times\_fourth\_eq\_octave} \\
\texttt{Hierarchy}
  & \texttt{extended\_ratio\_cost\_hierarchy} \\
  & \texttt{superparticular\_jcost} \\
  & \texttt{superparticular\_jcost\_decreasing} \\
  & \texttt{consonance\_is\_jcost} \\
\texttt{Fifth as bridge}
  & \texttt{twelve\_eight\_ratio\_is\_fifth} \\
  & \texttt{pythagorean\_comma\_positive} \\
  & \texttt{pythagorean\_comma\_small} \\
  & \texttt{chromaticBridgeRatio\_eq\_fifth} \\
  & \texttt{fifth\_is\_structural\_not\_lowest\_cost} \\
  & \texttt{twelve\_step\_bridge\_is\_low\_error\_witness} \\
\texttt{Valence}
  & \texttt{major\_skew\_gt\_minor\_skew} \\
  & \texttt{major\_third\_skew\_pos}, \texttt{minor\_third\_skew\_pos} \\
  & \texttt{ledgerSkew\_antisymmetric} \\
  & \texttt{ledgerSkew\_neg\_below\_one} \\
  & \texttt{skew\_increases\_with\_interval\_size} \\
  & \texttt{valence\_difference\_one\_twelfth} \\
\texttt{Conflict/coherence}
  & \texttt{chordConflict\_nonneg} \\
  & \texttt{chordCoherence\_pos}, \texttt{chordCoherence\_le\_one} \\
  & \texttt{lowerConflict\_higherCoherence} \\
\bottomrule
\end{tabular}
\end{table}

\noindent The full source code is available at
\texttt{github.com/jonwashburn/reality} in the
\texttt{IndisputableMonolith} directory.

% ====================================================================
\section{Discussion}
\label{sec:discussion}

\subsection{\texorpdfstring{What $J$-cost captures}{What J-cost captures}}

The $J$-cost function provides a single, parameter-free measure of
interval strain that is motivated by four transparent axioms (inversion
symmetry, unit normalisation, strict convexity, calibration) and
uniqueness-justified in the current formalization by a stronger explicit
functional-equation hypothesis bundle.  It
reproduces the qualitative observation that unison is maximally
consonant and the octave is the most ``costly'' standard interval, and
it provides exact rational values for every just-intonation ratio.  The
superparticular cost formula $J((n+1)/n)=1/(2n(n+1))$ gives a clean
parametric family that covers all of the traditional intervals.

\subsection{\texorpdfstring{Where $J$-cost diverges from tradition}{Where J-cost diverges from tradition}}

The most notable divergence is the ranking of the perfect fifth.
Traditional frameworks place the fifth immediately after the unison and
octave in consonance.  $J$-cost places it sixth out of seven standard
intervals, behind both thirds and the fourth.

This is not an error in the model; it is a structural prediction.  $J$
measures deviation from unity, and $3/2$ deviates more than $6/5$,
$5/4$, or $4/3$.  The fifth's cultural and structural prominence comes
instead from its bridge property (Section~\ref{sec:fifth}).  Whether
perceptual consonance tracks $J$-cost or roughness/template models is
an empirical question.

\subsection{The valence prediction}

The $\sigma$-function prediction developed here is comparative rather
than categorical: both the major and minor third have positive skew in
the current formalization, but the major third has more.  This yields a
parameter-free comparative account of major/minor brightness ordering
that is testable via affective ratings of isolated intervals.

\subsection{Limitations}

\begin{enumerate}
  \item \textbf{Just intonation only.}  The current framework applies
    to exact rational ratios.  Extension to equal temperament, where
    all intervals are irrational, requires a notion of ``nearest
    rational approximation cost'' that is not yet formalised.

  \item \textbf{Timbre is not modelled.}  $J$-cost applies to
    fundamental frequency ratios and does not account for the spectral
    content of real instruments, which affects perceived consonance.

  \item \textbf{Cultural context.}  Consonance perception has a
    well-documented cultural component \citep{mcdermott2016}.  $J$-cost
    makes structural predictions about cost that may or may not align
    with learned preference in every population.

  \item \textbf{Uniqueness hypotheses.}  The uniqueness theorem
    (Theorem~\ref{thm:unique}) requires the d'Alembert functional
    equation together with additional explicit regularity hypotheses.
    This package is natural in the Recognition Science framework but
    should be evaluated on its own merits in a purely music-theoretic
    context.
\end{enumerate}

\subsection{Future directions}

\paragraph{Non-Western tuning systems.}
Applying $J$-cost to the 22-\'{s}ruti system of Indian classical music,
the 17-tone Arabic system, or Indonesian gamelan scales would test
whether the cost hierarchy has cross-cultural predictive power.

\paragraph{Chord-level predictions.}
The conflict-coherence framework (Section~\ref{sec:conflict}) predicts
that among triads, the chord with the lowest aggregate $J$-cost should
produce the highest coherence.  This can be tested with EEG
inter-frequency coherence measurements.

\paragraph{Integration with roughness models.}
$J$-cost and roughness models may capture orthogonal aspects of
consonance perception.  A combined model---$J$-cost for structural
strain, roughness for sensory dissonance---could outperform either
alone.

% ====================================================================
\section{Conclusion}
\label{sec:conclusion}

We have shown that a single function $J(r)=\tfrac{1}{2}(r+r^{-1})-1$,
uniqueness-justified in the present formalization by an explicit bundle
of symmetry, normalization, convexity, calibration, continuity, and
functional-equation hypotheses, simultaneously
produces:

\begin{itemize}
  \item a complete, strict ordering of standard musical intervals,
  \item a structural explanation for the perfect fifth as the chromatic
    bridge ratio rather than the lowest-cost interval,
  \item a parameter-free comparative model of valence via the signed
    asymmetry $\sigma(r)=r-r^{-1}$,
  \item a monotone conflict-coherence framework for chords.
\end{itemize}

\noindent The framework makes specific predictions that differ from
traditional consonance theories, particularly regarding the relative
ranking of the perfect fifth.  The core theorem surface cited in the
paper has been machine-checked in a proof assistant.  We invite
empirical testing of the divergent predictions.

% ====================================================================
\section*{Acknowledgments}

The author thanks the Recognition Science community for foundational
discussions on cost-based physics, and the Lean~4 / Mathlib communities
for the proof assistant infrastructure that enabled formal verification
of these results.

% ====================================================================
\begin{thebibliography}{99}

\bibitem[Acz{\'e}l, 1966]{aczel1966}
Acz{\'e}l, J. (1966).
\textit{Lectures on Functional Equations and Their Applications}.
Academic Press.

\bibitem[Amiot, 2016]{amiot2016}
Amiot, E. (2016).
\textit{Music Through Fourier Space: Discrete Fourier Transform in
Music Theory}. Springer.

\bibitem[Benson, 2006]{benson2006}
Benson, D.~J. (2006).
\textit{Music: A Mathematical Offering}. Cambridge University Press.

\bibitem[Helmholtz, 1863]{helmholtz1863}
von Helmholtz, H. (1863).
\textit{Die Lehre von den Tonempfindungen als physiologische Grundlage
f\"ur die Theorie der Musik}. Braunschweig: Vieweg.

\bibitem[Mazzola, 2002]{mazzola2002}
Mazzola, G. (2002).
\textit{The Topos of Music: Geometric Logic of Concepts, Theory, and
Performance}. Birkh\"auser.

\bibitem[McDermott et al., 2016]{mcdermott2016}
McDermott, J.~H., Schultz, A.~F., Undurraga, E.~A., \& Godoy, R.~A.
(2016).
Indifference to dissonance in native Amazonians reveals cultural
variation in music perception.
\textit{Nature}, 535(7613), 547--550.

\bibitem[de Moura et al., 2021]{moura2021}
de Moura, L., Kong, S., Avigad, J., van Doorn, F., \& von Raumer, J.
(2021).
The Lean 4 theorem prover and programming language.
\textit{Automated Deduction---CADE 28}, LNAI 12699, 625--635.

\bibitem[Parncutt \& Hair, 2011]{parncutt2011}
Parncutt, R. \& Hair, G. (2011).
Consonance and dissonance in music theory and psychology: disentangling
roughness, harmonicity, and familiarity.
\textit{Journal of the Acoustical Society of America}, 130(2), 776--788.

\bibitem[Plomp \& Levelt, 1965]{plomp1965}
Plomp, R. \& Levelt, W.~J.~M. (1965).
Tonal consonance and critical bandwidth.
\textit{Journal of the Acoustical Society of America}, 38(4), 548--560.

\bibitem[Washburn, 2025]{washburn2025axioms}
Washburn, J. (2025).
Recognition Science: a cost-first foundation for physics.
\textit{Axioms}, 14(x).

\bibitem[Washburn, 2026]{washburn2026envelope}
Washburn, J. (2026).
Precision envelope targeting of neural oscillation bands: exact
rhythmic formulas, a conflict-coherence framework, and a unique carrier
theorem for the theta band.
\textit{Manuscript submitted for publication}.

\end{thebibliography}

\end{document}
